\section{관련 연구}
\label{sec:related}

\subsection{추론 보강형 E2E/VLA 자율주행}

DriveLM은 주행 장면을 그래프 기반 시각 질의응답으로 구조화하고 \cite{sima2024drivelm}, DriveVLM은 VLM의 장면 이해와 계획을 결합한다 \cite{tian2024drivevlm}. Alpamayo-R1은 CoC 추론과 행동/궤적 예측을 연결하여 긴 꼬리 상황에서의 일반화를 목표로 한다 \cite{wang2025alpamayo}. 한편 자율주행 VLM의 신뢰성 연구는 장면 이해와 행동 추론의 평가 지표 및 데이터 품질 문제를 지적한다 \cite{xie2025drivebench}. 본 연구는 객체 환각이나 CoC의 일반적 사실 오류를 다시 검출하기보다, 동일한 행동 목표를 달성하는 후보 안에서 명시된 허용조건을 학습하고 적용할 수 있는지를 다룬다.

\subsection{주행 안전 제약과 검증}

자율주행 안전 연구는 안전거리, 책임 민감 안전 모델, 도달 가능성, 실행시간 감시기(runtime monitor)와 차단장치(shield)처럼 실행 단계에서 제약을 검사하는 방법을 발전시켜 왔다 \cite{shalev2017rss,katz2019marabou}. 시나리오 명세 및 폐루프 시험은 보지 않은 조합에서 학습 정책의 취약성을 발견하는 데 사용된다 \cite{fremont2019scenic,zhang2023cat}. 이러한 방법은 정밀한 제약 집행에 적합하지만 자연어 추론과 실행 계약 사이의 인터페이스는 별도 문제다.

\subsection{자연어와 정형 명세의 역할}

자연어는 사전학습 VLM이 처리하기 쉬우며 사람에게 편리하지만, 부정 범위, 단위 및 경계 연산자가 모호할 수 있다. 정형 명세는 정확한 판정과 실행 단계의 강제 적용에 적합하지만 기반 VLM이 임의 논리 문법을 즉시 실행할 것이라고 가정할 수 없다. 본 논문은 두 표현을 경쟁 관계로 두지 않는다. 자연어 \sccoc{}를 학습 인터페이스로 평가하고, 수치 판정은 모델 외부의 결정적 검증기로 계산한다. 자연어에서 구조화 계약(typed contract)으로의 변환은 후속 연구로 남긴다.

\subsection{차별점}

본 연구는 실제 사고 궤적을 필요로 하지 않는다. 정적 과제에서는 모든 후보가 CoC 목표를 달성하되 가드 준수 여부와 진행도가 다르게 구성하고, 시간 과제에서는 계속 정지하는 교착 대조군을 추가한다. 또한 동일 장면의 계약만 바꾸는 쌍대 계약과 제약--정답 대응 교란 통제군을 사용하여 위반 감소와 목표 완료를 함께 평가한다.
